\documentclass[a4paper,12pt]{article}
\usepackage[utf8]{inputenc}
\usepackage[T1]{fontenc}
\usepackage{amsmath}
\usepackage{booktabs}
\usepackage{graphicx}
\usepackage{geometry}
\usepackage{xcolor}
\usepackage{listings}
\usepackage{hyperref}

\geometry{margin=2.5cm}
\hypersetup{colorlinks=true, linkcolor=blue!50!black, urlcolor=blue!50!black}

\renewcommand{\tablename}{Tabela}
\renewcommand{\lstlistingname}{Código}



% ---------- Configuração do realce de sintaxe (C) ----------

\definecolor{codebg}{RGB}{247,247,247}
\definecolor{codekw}{RGB}{0,0,180}
\definecolor{codecomment}{RGB}{0,128,0}
\definecolor{codestring}{RGB}{163,21,21}
\definecolor{framecolor}{RGB}{200,200,200}

\lstdefinestyle{cstyle}{
    language=C,
    backgroundcolor=\color{codebg},
    basicstyle=\ttfamily\footnotesize,
    keywordstyle=\color{codekw}\bfseries,
    commentstyle=\color{codecomment}\itshape,
    stringstyle=\color{codestring},
    numberstyle=\tiny\color{gray},
    numbers=left,
    numbersep=6pt,
    frame=single,
    rulecolor=\color{framecolor},
    breaklines=true,
    breakatwhitespace=false,
    showstringspaces=false,
    tabsize=4,
    captionpos=b,
    columns=flexible,
    morekeywords={size_t},
    extendedchars=true,
    literate=
        {á}{{\'a}}1 {à}{{\`a}}1 {â}{{\^a}}1 {ã}{{\~a}}1
        {é}{{\'e}}1 {ê}{{\^e}}1 {í}{{\'i}}1
        {ó}{{\'o}}1 {ô}{{\^o}}1 {õ}{{\~o}}1
        {ú}{{\'u}}1 {ü}{{\"u}}1 {ç}{{\c c}}1
        {Á}{{\'A}}1 {À}{{\`A}}1 {Â}{{\^A}}1 {Ã}{{\~A}}1
        {É}{{\'E}}1 {Ê}{{\^E}}1 {Í}{{\'I}}1
        {Ó}{{\'O}}1 {Ô}{{\^O}}1 {Õ}{{\~O}}1
        {Ú}{{\'U}}1 {Ü}{{\"U}}1 {Ç}{{\c C}}1
}



\title{Tarefa 4: Aplicações limitadas por memória ou CPU}
\author{Israel Soares de Castro Filho \\ Matrícula: 20250070780}
\date{}

\begin{document}
\maketitle

\section{Introdução}

Este relatório apresenta a implementação e a análise de dois programas paralelos
em C, utilizando a biblioteca OpenMP, com o objetivo de contrastar dois padrões
distintos de gargalo computacional: aplicações \emph{memory-bound} (limitadas
pela banda de memória) e aplicações \emph{compute-bound} (limitadas pela
capacidade de processamento aritmético da CPU).

O primeiro programa realiza uma soma simples de vetores, operação de baixíssima
intensidade aritmética e alto volume de dados, representativa de cargas
limitadas por memória. O segundo programa aplica uma sequência de operações
matemáticas (\texttt{sin}, \texttt{cos}, \texttt{sqrt}) repetidas diversas vezes
sobre um vetor bem menor, representando uma carga limitada pela capacidade de
cálculo da CPU. Ambos foram paralelizados com a diretiva
\texttt{\#pragma omp parallel for} e medidos variando o número de threads
(1, 2, 4, 8 e 16), a fim de observar como cada tipo de aplicação responde ao
aumento de paralelismo e, em particular, ao paralelismo de hardware
disponível acima do número de núcleos físicos.

A máquina utilizada nos testes possui \textbf{14 núcleos físicos e 20
processadores lógicos}. Esse número não é explicado por uma arquitetura
homogênea com Hyper-Threading simples (o que exigiria um número par de
threads por núcleo), mas sim por uma \textbf{arquitetura híbrida} (como as
CPUs Intel de 12\textsuperscript{a}/13\textsuperscript{a} geração, no estilo
\emph{Performance-cores} + \emph{Efficiency-cores}): considerando que apenas
os núcleos de desempenho (P-cores) possuem Hyper-Threading (2 threads cada) e
os núcleos de eficiência (E-cores) possuem apenas 1 thread cada, a única
combinação que resulta em 14 núcleos e 20 threads é
$6~\text{P-cores} \times 2 + 8~\text{E-cores} \times 1 = 12 + 8 = 20$ threads,
com $6 + 8 = 14$ núcleos físicos. Essa heterogeneidade de núcleos é essencial
para interpretar corretamente os resultados de escalabilidade apresentados
neste relatório.

\section{Metodologia}

Os dois programas foram escritos em C e paralelizados com OpenMP utilizando
exclusivamente a diretiva \texttt{\#pragma omp parallel for schedule(static)},
sem uso de \texttt{reduction}, já que cada iteração escreve em uma posição de
vetor independente, não havendo condição de corrida.

\begin{itemize}
    \item \textbf{Programa 1 (\texttt{memory\_bound.c})}: aloca três vetores de
    \texttt{double} de tamanho $n = 60{.}000{.}000$ (aproximadamente 480~MB por
    vetor) e executa a operação elemento a elemento $c[i] = a[i] + b[i]$.
    A intensidade aritmética é mínima (uma soma para dois acessos de leitura e
    um de escrita), de modo que o tempo de execução é dominado pelo tráfego
    entre a CPU e a memória principal.

    \item \textbf{Programa 2 (\texttt{compute\_bound.c})}: aloca vetores bem
    menores ($n = 2{.}000{.}000$ elementos) e, para cada elemento, executa um
    laço interno de 400 iterações contendo chamadas a \texttt{sin}, \texttt{cos}
    e \texttt{sqrt}. A intensidade aritmética é alta e o volume de dados é
    pequeno o suficiente para caber majoritariamente em cache, de modo que o
    tempo de execução é dominado pela capacidade de cálculo das unidades de
    ponto flutuante.
\end{itemize}

Em ambos os programas, os vetores de entrada são inicializados dentro de uma
região paralela com o mesmo \texttt{schedule(static)} usado no laço medido,
garantindo a política de \emph{first-touch}. 
A medição de tempo utiliza \texttt{omp\_get\_wtime()}
(tempo de parede), e cada configuração é executada 5 vezes, sendo reportado o
melhor tempo (\texttt{best}) entre as repetições.

Os programas foram compilados com:
\begin{lstlisting}[style=cstyle, language=bash, numbers=none]
gcc -fopenmp -O2 memory_bound.c  -o memory_bound  -lm
gcc -fopenmp -O2 compute_bound.c -o compute_bound -lm
\end{lstlisting}

e executados para $1$, $2$, $4$, $8$ e $16$ threads, mantendo os demais
parâmetros fixos (tamanho do vetor, número de iterações internas e número de
repetições).

\subsection{Identificação do hardware}

Antes de interpretar os resultados, foi executado um pequeno programa de
diagnóstico (Apêndice~\ref{sec:cpu_info}) que consulta
\texttt{omp\_get\_num\_procs()} e \texttt{omp\_get\_max\_threads()} para
determinar quantos processadores lógicos o OpenMP enxerga por padrão. A
saída obtida foi:

\begin{lstlisting}[style=cstyle, language=bash, numbers=none]
Number of available cores: 20
Max OpenMP threads: 20
Running with 20 active threads
\end{lstlisting}

Combinando essa saída com a informação fornecida pelo fabricante/sistema
operacional de que a máquina possui \textbf{14 núcleos físicos}, conclui-se
que se trata de uma CPU com arquitetura híbrida (núcleos de desempenho e de
eficiência), e não de uma CPU homogênea com Hyper-Threading simples em todos
os núcleos. Essa distinção é retomada na Seção~\ref{sec:analise} para
explicar o comportamento do programa compute-bound.

\section{Resultados}

As Tabelas~\ref{tab:memoria} e \ref{tab:compute} apresentam o melhor tempo de
execução (em segundos) obtido para cada número de threads, junto com o
\emph{speedup} (em relação à execução com 1 thread) e a eficiência paralela
($\text{speedup}/\text{threads}$).

\begin{table}[h]
    \centering
    \caption{Resultados do programa memory-bound ($n = 60.000.000$ elementos)}
    \label{tab:memoria}
    \begin{tabular}{ccccc}
        \toprule
        Threads & Tempo (s) & Speedup & Eficiência (\%) \\
        \midrule
        1  & 0.047 & 1.00$\times$ & 100.0 \\
        2  & 0.032 & 1.47$\times$ & 73.4  \\
        4  & 0.021 & 2.24$\times$ & 56.0  \\
        8  & 0.021 & 2.24$\times$ & 28.0  \\
        16 & 0.020 & 2.35$\times$ & 14.7  \\
        \bottomrule
    \end{tabular}
\end{table}

\begin{table}[h]
    \centering
    \caption{Resultados do programa compute-bound ($n = 2.000.000$ elementos)}
    \label{tab:compute}
    \begin{tabular}{ccccc}
        \toprule
        Threads & Tempo (s) & Speedup & Eficiência (\%) \\
        \midrule
        1  & 14.381 & 1.00$\times$  & 100.0 \\
        2  & 7.157  & 2.01$\times$  & 100.5 \\
        4  & 3.600  & 4.00$\times$  & 99.9  \\
        8  & 1.964  & 7.32$\times$  & 91.5  \\
        16 & 1.167  & 12.32$\times$ & 77.0  \\
        \bottomrule
    \end{tabular}
\end{table}

Observa-se um comportamento claramente distinto entre os dois programas:
o programa compute-bound apresenta \emph{speedup} praticamente linear até 8
threads, ao passo que o programa memory-bound satura rapidamente, já
apresentando ganhos marginais a partir de 4 threads e nenhum ganho relevante
entre 4 e 8 threads.

\section{Análise}
\label{sec:analise}

\subsection{Programa memory-bound}

A operação $c[i] = a[i] + b[i]$ tem intensidade aritmética muito baixa
(duas leituras e uma escrita para cada soma), então o gargalo não é a CPU, e
sim a banda de memória, um recurso compartilhado por todos os núcleos. Por
isso a eficiência já cai de 100\% para 73{,}4\% com apenas 2 threads e a
banda satura por volta de 4 threads: o tempo com 4 e com 8 threads é
praticamente o mesmo (0{,}021~s). O pequeno ganho entre 8 e 16 threads
(2{,}24$\times$ para 2{,}35$\times$) é irrelevante frente aos tempos curtos
medidos (dezenas de milissegundos) e pode ser atribuído a ruído.

\subsection{Programa compute-bound}

A máquina de testes tem 14 núcleos físicos e 20 threads lógicas, divididos
em 6 P-cores (rápidos, com Hyper-Threading) e 8 E-cores (mais lentos, sem
Hyper-Threading). Isso explica a curva de eficiência em duas etapas:
 
\textbf{Até 8 threads}, ainda há núcleos físicos livres, então a eficiência
se mantém alta (100\% até 4 threads, 91{,}5\% em 8). A pequena queda em 8
threads já sugere que parte delas foi parar nos E-cores, mais lentos que os
P-cores.
 
\textbf{Em 16 threads}, a eficiência cai mais (77{,}0\%) por dois motivos
somados: mais threads rodando nos E-cores mais lentos, e o fato de 16
ultrapassar os 14 núcleos físicos, obrigando algumas threads a dividir um
mesmo P-core via Hyper-Threading — disputando a mesma FPU e a mesma cache.
Ainda assim o speedup continua crescendo (7{,}32$\times$ para
12{,}32$\times$), pois os E-cores ainda somam capacidade real de cálculo, só
que bem abaixo do que se teria com 16 núcleos físicos homogêneos.

\subsection{Comparação entre os dois padrões}

Os resultados confirmam a hipótese da atividade: threads de hardware
(Hyper-Threading) ajudam mais em cargas memory-bound do que em cargas
compute-bound.
 
\begin{itemize}
    \item \textbf{Memory-bound}: o Hyper-Threading poderia ajudar
    preenchendo o tempo ocioso de uma thread esperando dados da memória, mas
    aqui o gargalo de banda já aparece bem antes (4 threads, longe dos 14
    núcleos físicos), então esse benefício nem chega a ser testado.
    \item \textbf{Compute-bound}: como as unidades de execução já ficam
    ocupadas o tempo todo, o Hyper-Threading só disputa recursos em vez de
    ajudar, e a presença de E-cores mais lentos reduz ainda mais a
    eficiência conforme mais threads são criadas.
\end{itemize}
 
Em resumo: o desempenho melhora quase linearmente enquanto há núcleos
físicos livres; estabiliza quando a memória satura; e a eficiência cai
progressivamente quando o paralelismo passa a depender de núcleos mais
lentos (E-cores) ou de Hyper-Threading em núcleos já ocupados.

\section{Conclusão}

A implementação e os experimentos realizados permitiram observar de forma
prática a diferença de comportamento entre aplicações limitadas por memória e
aplicações limitadas por CPU quando paralelizadas com OpenMP. O programa
memory-bound atingiu seu limite de escalabilidade rapidamente (por volta de 4
threads), refletindo a saturação da banda de memória, um recurso
compartilhado entre todos os núcleos, bem abaixo dos 14 núcleos físicos
disponíveis na máquina de testes. Já o programa compute-bound escalou de
forma quase linear enquanto pôde ocupar núcleos físicos distintos, e teve sua
eficiência reduzida gradualmente à medida que (i) mais threads passaram a
ser executadas nos núcleos de eficiência (E-cores), mais lentos, e (ii),
acima de 14 threads, algumas delas passaram a compartilhar um mesmo núcleo
de desempenho (P-core) via Hyper-Threading.
 
Um resultado adicional relevante deste trabalho foi a constatação, por meio
do programa de diagnóstico (\texttt{omp\_get\_num\_procs()}), de que a
máquina utilizada não possui uma arquitetura homogênea de núcleos, mas sim
uma arquitetura híbrida (P-cores + E-cores), com 14 núcleos físicos e 20
threads lógicas. Isso demonstra, na prática, que a simples contagem de
``processadores lógicos'' reportada pelo sistema operacional ou pelo OpenMP
não é suficiente para prever o comportamento de escalabilidade de uma
aplicação: é necessário também saber quantos desses processadores são
núcleos físicos independentes, quantos compartilham recursos via
Hyper-Threading, e se todos os núcleos físicos têm o mesmo desempenho.
Esses resultados reforçam a importância de identificar tanto o tipo de
gargalo de uma aplicação (CPU ou memória) quanto a topologia real do
hardware disponível antes de decidir a estratégia de paralelização e o
número de threads mais adequado.

\newpage

\appendix
\section{Código-Fonte}

\subsection{memory\_bound.c}
\lstinputlisting[style=cstyle]{../src/memory_bound.c}

\subsection{compute\_bound.c}
\lstinputlisting[style=cstyle]{../src/compute_bound.c}

\subsection{cpu\_info.c --- programa de diagnóstico de hardware}
\label{sec:cpu_info}
\lstinputlisting[style=cstyle]{../src/cpu_info.c}

\end{document}